\anexo{Glosario preliminar}

\begin{description}[leftmargin=0cm, labelsep=0.5cm]
    \item[\textbf{Bovino de cr\'ia:}] Animal perteneciente a un sistema productivo orientado a la reproducci\'on y posterior venta.
    \item[\textbf{Caravana:}] Identificador colocado en la oreja del animal, utilizado para reconocerlo de forma individual dentro del rodeo.
    \item[\textbf{Lote:}] Grupo de animales organizado seg\'un criterios productivos, sanitarios, etarios o de manejo.
    \item[\textbf{Tacto:}] Pr\'actica veterinaria utilizada para evaluar el estado reproductivo de una hembra bovina.
    \item[\textbf{Sangrado:}] Extracci\'on de muestra de sangre para an\'alisis sanitario o reproductivo.
    \item[\textbf{Rodeo:}] Conjunto de animales de un establecimiento ganadero.
    \item[\textbf{Ganader\'ia de precisi\'on:}] Uso de tecnolog\'ias de monitoreo, registro y an\'alisis de datos para mejorar la gesti\'on ganadera.
\end{description}

Este glosario deber\'a ampliarse con la terminolog\'ia definitiva utilizada en la tesis.
